\documentclass[11pt]{article}

\usepackage{amsmath, amssymb, amsfonts}
\usepackage{geometry}
\geometry{margin=1in}

\title{Supplementary Information: Information Geometry and Score Fluctuation Consistency}
\author{}
\date{}

\begin{document}

\maketitle

\section*{1. Log-Likelihood and Score Definitions}

Consider a trajectory of $T$ sequential measurement intervals.
The total log-likelihood is defined as
\begin{equation}
\ell(\theta) 
= \log L(\theta)
= \sum_{t=1}^T \ell_t(\theta),
\end{equation}
where $\ell_t(\theta)$ denotes the log-likelihood contribution of interval $t$.

The total score is
\begin{equation}
S(\theta) = \nabla_\theta \ell(\theta)
= \sum_{t=1}^T s_t(\theta),
\end{equation}
with per-interval score
\begin{equation}
s_t(\theta) = \nabla_\theta \ell_t(\theta).
\end{equation}

Throughout, expectations and covariances are taken over independent realizations of full trajectories
generated under parameter $\theta$.

\section*{2. Gaussian Predictive Model and Per-Interval Log-Likelihood}

At each interval $t$, the predictive likelihood is approximated as Gaussian,
\begin{equation}
y_t \,\big|\, \mathcal{F}_{t-1} \;\sim\; \mathcal{N}\!\left(\mu_t(\theta), \sigma_t^2(\theta)\right),
\end{equation}
where $\mu_t(\theta)$ and $\sigma_t^2(\theta)$ are produced recursively by the inference algorithm
(e.g.\ MacroIR / MacroMR) given past information $\mathcal{F}_{t-1}$.

The corresponding per-interval log-likelihood is
\begin{equation}
\ell_t(\theta)
=
-\frac{1}{2}\log\!\left(2\pi\sigma_t^2(\theta)\right)
-\frac{\left(y_t-\mu_t(\theta)\right)^2}{2\sigma_t^2(\theta)}.
\end{equation}

\section*{3. Gaussian Fisher Information (GFI)}

We define the Gaussian Fisher Information matrix for the \emph{total} log-likelihood as
\begin{equation}
\mathbf H(\theta)
:=
- \mathbb{E}\left[ \nabla_\theta^2 \ell(\theta) \right].
\end{equation}

In practice, we compute $\mathbf H(\theta)$ by summing the per-interval Fisher information
implied by the Gaussian predictive model. Let
\begin{equation}
\mathbf g_t(\theta) := \nabla_\theta \mu_t(\theta) \in \mathbb{R}^d,
\qquad
\mathbf h_t(\theta) := \nabla_\theta \sigma_t^2(\theta) \in \mathbb{R}^d,
\end{equation}
with scalar predictive variance $\sigma_t^2(\theta) > 0$.

\subsection*{3.1 Per-interval Gaussian Fisher information}

For a univariate normal $\mathcal{N}(\mu_t(\theta),\sigma_t^2(\theta))$ with $\sigma_t^2$ scalar,
the per-interval Fisher information is
\begin{equation}
\mathbf I_t^{\mathrm{G}}(\theta)
=
\frac{1}{\sigma_t^2(\theta)}\,\mathbf g_t(\theta)\mathbf g_t(\theta)^\top
+
\frac{1}{2\,\sigma_t^4(\theta)}\,\mathbf h_t(\theta)\mathbf h_t(\theta)^\top.
\label{eq:gaussian_fisher_interval}
\end{equation}

This expression follows from the standard Fisher information of a normal distribution,
combined with the chain rule for $\mu_t(\theta)$ and $\sigma_t^2(\theta)$.

\subsection*{3.2 Total Gaussian Fisher information}

The total Gaussian Fisher Information is then obtained by summation (or equivalently, by the curvature of the total log-likelihood):
\begin{equation}
\mathbf H(\theta)
\approx
\sum_{t=1}^T \mathbb{E}\!\left[\mathbf I_t^{\mathrm{G}}(\theta)\right].
\label{eq:gaussian_fisher_total}
\end{equation}

In Monte Carlo evaluation (many independent simulated trajectories), the expectation in
Eq.~\eqref{eq:gaussian_fisher_total} is estimated by averaging $\mathbf I_t^{\mathrm{G}}(\theta)$ across realizations
at each $t$, then summing over $t$.

\section*{4. Score Covariance Matrix (SCM)}

We define the Score Covariance Matrix as the covariance of the \emph{total} score,
\begin{equation}
\mathbf J(\theta)
:=
\mathrm{Cov}\!\left( S(\theta) \right)
=
\mathbb{E}\!\left[
\left(S(\theta) - \mathbb{E}[S(\theta)]\right)
\left(S(\theta) - \mathbb{E}[S(\theta)]\right)^{\!\top}
\right].
\end{equation}

Unlike $\mathbf H$, which is curvature-implied by the Gaussian predictive model through
Eq.~\eqref{eq:gaussian_fisher_interval}, $\mathbf J$ is fluctuation-implied: it is determined by realized variability of the
total score across independent trajectories, including any temporal correlations between intervals.

\section*{5. Information Distortion Matrix (IDM)}

To compare curvature-implied and fluctuation-implied information geometries, we define the
Information Distortion Matrix
\begin{equation}
\mathbf C(\theta)
=
\mathbf H(\theta)^{-1/2}
\mathbf J(\theta)
\mathbf H(\theta)^{-1/2}.
\end{equation}

Here $\mathbf H^{-1/2}$ denotes the inverse square root of $\mathbf H$.
In computations, we obtain $\mathbf H^{-1/2}$ via a Cholesky factorization
$\mathbf H = \mathbf L\mathbf L^\top$ and triangular solves, without forming explicit matrix inverses.

The matrix $\mathbf C$ is symmetric positive definite and dimensionless.

\section*{6. Information Consistency}

If the Gaussian predictive model fully captures the conditional dynamics relevant to the likelihood increments,
then curvature and fluctuations are consistent at the true parameter $\theta_0$:
\begin{equation}
\mathbf J(\theta_0) = \mathbf H(\theta_0),
\qquad
\mathbf C(\theta_0) = \mathbf I.
\end{equation}

In this regime, the predictive model acts as an effective whitening transformation for the information channel:
the total score fluctuations match the curvature predicted by the Gaussian Fisher geometry.

\section*{7. Relation to Estimator Covariance}

The covariance of the maximum likelihood estimator is given asymptotically by the (misspecification-robust) sandwich form,
\begin{equation}
\mathrm{Cov}(\hat{\theta})
=
\mathbf H^{-1} \mathbf J \mathbf H^{-1}.
\end{equation}

Thus $\mathbf C$ characterizes the intrinsic information distortion that separates the naive Fisher covariance $\mathbf H^{-1}$
from the fluctuation-corrected covariance.

\section*{8. Spectral Interpretation}

Let $\lambda_i$ denote the eigenvalues of $\mathbf C$.
Then:
\begin{itemize}
\item $\lambda_i = 1$ indicates perfect agreement between curvature and fluctuation geometry along an eigen-direction.
\item $\lambda_i > 1$ indicates variance inflation relative to the Gaussian Fisher prediction.
\item $\lambda_i < 1$ indicates deflation (overestimation of curvature).
\end{itemize}

The determinant $\det(\mathbf C)$ quantifies global information-volume distortion, while $\mathrm{tr}(\mathbf C)/d$ measures
average directional distortion.

\section*{9. Interpretation}

The Gaussian Fisher Information $\mathbf H$ encodes the local quadratic geometry implied by the Gaussian predictive approximation
through Eqs.~\eqref{eq:gaussian_fisher_interval}--\eqref{eq:gaussian_fisher_total}. The Score Covariance Matrix $\mathbf J$ encodes
the geometry implied by realized score fluctuations across trajectories. The Information Distortion Matrix $\mathbf C$ measures the
mismatch between these two geometries. When $\mathbf C=\mathbf I$, the Gaussian predictive model has successfully whitened the
information channel of the trajectory, in the sense that curvature-implied and fluctuation-implied geometries agree.

\end{document}
